% ============================================================
%  CHAPTER 2: LITERATURE SURVEY
% ============================================================
\chapter{Literature Survey}

This chapter reviews the existing research and techniques relevant to the
problem addressed in this project. Section~2.1 covers traditional routing and
vehicle scheduling methods. Section~2.2 discusses the application of machine
learning to logistics and supply chain forecasting. Section~2.3 provides the
theoretical foundations of reinforcement learning, including Markov Decision
Processes, Q-Learning, Deep Q-Networks, and Proximal Policy Optimization, in
the level of detail required to understand the design decisions made in this
project.

% ================================================================
\section{Traditional Routing and Vehicle Scheduling Methods}

% ----------------------------------------------------------------
\subsection{Dijkstra's Algorithm and A* Search}

Graph-based shortest-path algorithms have been the foundation of computational
route planning since the late 1950s. Dijkstra's algorithm, introduced in 1959
\cite{dijkstra1959}, computes the shortest path from a single source node to all
other nodes in a weighted graph. It is guaranteed to find the optimal path when
all edge weights are non-negative and deterministic. For logistics applications,
the road network is modeled as a weighted directed graph where nodes represent
intersections or depots and edge weights represent travel time or distance.

The A* algorithm, introduced by Hart, Nilsson, and Raphael \cite{hart1968},
extends Dijkstra's approach by incorporating a heuristic function that estimates
the remaining distance to the destination. This makes A* substantially faster in
practice for point-to-point routing. In the baseline configuration of this
project, a modified A* algorithm is used as the base route planner. However, both
Dijkstra's algorithm and A* share the fundamental limitation that their cost
functions are fixed at the time of planning. They cannot respond to events that
occur while the truck is already in transit, such as a road accident that blocks
a segment, or a sudden increase in traffic density in a particular zone. This
static nature makes them unsuitable as a standalone solution for dynamic fleet
management.

% ----------------------------------------------------------------
\subsection{The Vehicle Routing Problem and Its Variants}

The Vehicle Routing Problem (VRP), formalized by Dantzig and Ramser in 1959 and
surveyed comprehensively by Toth and Vigo \cite{toth2002}, seeks to find a set of
routes for a fleet of vehicles that minimizes total travel cost while satisfying
the delivery demand of a set of customers. The VRP is a generalization of the
Traveling Salesman Problem and is known to be NP-hard in the general case.

Several variants of the VRP are directly relevant to the supply chain routing
problem. The Capacitated VRP (CVRP) adds vehicle load capacity constraints. The
VRP with Time Windows (VRPTW) adds delivery time deadlines, which is analogous
to the cargo RSL deadline in this project: a delivery that arrives after the RSL
expires is effectively a failed delivery. Despite extensive research, VRP solvers
face two key limitations in the context of this project. First, they are
computationally expensive and do not scale to real-time dynamic rerouting for
large fleets. Second, they assume a fixed set of delivery requests known at the
start, whereas in a real supply chain orders arrive continuously and conditions
change mid-route.

% ----------------------------------------------------------------
\subsection{Limitations of Traditional Approaches}

Traditional routing methods share several fundamental limitations that motivate
the use of reinforcement learning in this project. Their cost functions are
defined at planning time and cannot adapt once the truck is en route. They
optimize a single objective, typically distance or time, and cannot dynamically
balance competing objectives such as fuel efficiency and cargo freshness. They do
not learn from historical patterns: a truck that repeatedly passes through a
congested zone at rush hour will receive no benefit from that experience in the
next planning cycle. These limitations collectively motivate the use of an
adaptive, learning-based approach.

% ================================================================
\section{Machine Learning in Supply Chain Management}

% ----------------------------------------------------------------
\subsection{Regression Models for ETA Forecasting}

Supervised machine learning methods have been applied to demand forecasting and
travel time estimation in logistics for over a decade. Ridge Regression
\cite{hoerl1970}, a regularized linear model that penalizes large coefficients
using an L2 norm, is well-suited to settings where input features are correlated
and the number of training samples grows incrementally over time. It can be
updated in an online fashion using partial fitting, making it computationally
lightweight and suitable for deployment on edge hardware without requiring a GPU.

In this project, Ridge Regression is used as the backbone of the ETA Forecaster.
The choice of a linear model over a deep neural network is deliberate: the ETA
Forecaster must update itself after every single road segment traversal, which
happens hundreds of times per simulated day. A deep neural network would require
batched gradient descent and would be far too slow for this use case. The Ridge
Regression model, by contrast, can update incrementally in milliseconds.

% ----------------------------------------------------------------
\subsection{Deep Learning for Sequential Forecasting}

Recurrent Neural Networks (RNNs) and Long Short-Term Memory (LSTM) networks have
demonstrated strong performance on sequential logistics data such as hourly demand
series and historical route time data \cite{sutton2018}. These models can capture
long-range temporal dependencies that linear models cannot. However, they require
substantial amounts of training data, fixed-size input sequences, and batch
training procedures that make them poorly suited for online, per-sample updates.

For the ETA Forecasting component of this project, the computational constraints
of online learning rule out LSTM-based approaches. The Ridge Regression model
achieves adequate accuracy for the purpose of route cost estimation while
remaining practical for real-time deployment.

% ----------------------------------------------------------------
\subsection{Advantages and Limitations of ML Approaches}

Machine learning models offer significant advantages over traditional statistical
methods in their ability to capture non-linear relationships and adapt to changing
data distributions. However, they remain fundamentally predictive: they forecast
what will happen but do not autonomously make decisions about what to do. A model
that accurately predicts delivery times cannot itself decide which route to take,
which truck to dispatch, or when to override a plan due to an emergency. These
decision-making requirements are precisely what reinforcement learning is designed
to address.

% ================================================================
\section{Reinforcement Learning in Logistics}

% ----------------------------------------------------------------
\subsection{Markov Decision Processes}

Reinforcement learning (RL) is a branch of machine learning in which an agent
learns to make sequential decisions by interacting with an environment and
receiving scalar reward signals \cite{sutton2018}. The theoretical foundation of
most RL algorithms is the Markov Decision Process (MDP), defined as a tuple
$(S, A, P, R, \gamma)$, where $S$ is the set of possible states, $A$ is the set
of actions, $P(s' | s, a)$ is the transition probability from state $s$ to state
$s'$ under action $a$, $R(s, a)$ is the reward received, and $\gamma \in [0, 1)$
is a discount factor that determines how much the agent values future rewards
relative to immediate ones.

The objective of the RL agent is to learn a policy $\pi(a | s)$ that maximizes
the expected cumulative discounted reward:

\begin{equation}
    J(\pi) = \mathbb{E}_{\pi} \left[ \sum_{t=0}^{\infty} \gamma^t R(s_t, a_t) \right]
\end{equation}

In the context of this project, the state $s$ is a 25-dimensional vector
describing the current condition of the truck and its environment (fuel level,
cargo RSL, traffic density, time of day, zone, and so on). The action $a$ is the
routing strategy selected by the Central PPO Brain: 0 for Fuel-Priority, 1 for
RSL-Priority, or 2 for Speed-Priority. The reward $R$ combines a delivery
completion bonus with penalties for cargo decay, fuel consumption, and retailer
stockouts.

% ----------------------------------------------------------------
\subsection{Q-Learning and Deep Q-Networks}

Q-Learning is a model-free, off-policy value-based RL algorithm introduced by
Watkins in 1989. It maintains a Q-function $Q(s, a)$ representing the expected
cumulative reward of taking action $a$ in state $s$ and then following the optimal
policy thereafter. The Q-function is updated after each step using the Bellman
equation:

\begin{equation}
    Q(s, a) \leftarrow Q(s, a) + \alpha \left[ r + \gamma \max_{a'} Q(s', a') - Q(s, a) \right]
\end{equation}

where $\alpha$ is the learning rate, $r$ is the immediate reward, and $s'$ is the
next state. In its tabular form, Q-Learning stores the Q-values in a table indexed
by all possible state-action pairs. This becomes impractical when the state space
is continuous or high-dimensional.

Deep Q-Networks (DQN), introduced by Mnih et al.\ \cite{mnih2015}, resolve this
limitation by replacing the Q-table with a deep neural network that approximates
the Q-function across continuous state spaces. Key stabilization techniques
introduced in DQN include experience replay (storing past transitions in a buffer
and sampling mini-batches for training) and the use of a separate target network
that is updated less frequently than the main network, which prevents oscillation
during training.

In this project, the Edge Q-Learning Agent uses the tabular form of Q-Learning.
The emergency state space (local hazard severity and number of available alternate
routes) is small and discrete, making a tabular approach practical. The Edge Agent
updates its Q-table online after each emergency event, allowing it to learn which
alternate routes are reliable in specific zones and time periods.

% ----------------------------------------------------------------
\subsection{Proximal Policy Optimization}

Proximal Policy Optimization (PPO), introduced by Schulman et al.\ \cite{schulman2017},
is a policy-gradient reinforcement learning algorithm that belongs to the
actor-critic family. Unlike value-based methods such as DQN, which learn a Q-value
function and derive the policy implicitly, PPO directly learns a parameterized
policy $\pi_\theta(a | s)$.

The key innovation of PPO is the clipped surrogate objective function, which
prevents the policy from changing too rapidly in a single update step:

\begin{equation}
    L^{CLIP}(\theta) = \mathbb{E}_t \left[ \min \left( r_t(\theta) \hat{A}_t, \;
    \text{clip}(r_t(\theta), 1 - \epsilon, 1 + \epsilon) \hat{A}_t \right) \right]
\end{equation}

where $r_t(\theta) = \frac{\pi_\theta(a_t | s_t)}{\pi_{\theta_{old}}(a_t | s_t)}$
is the probability ratio between the new and old policies, $\hat{A}_t$ is the
estimated advantage function, and $\epsilon$ is the clipping threshold (typically
0.2). The clip operation ensures that the policy update is bounded, preventing
large destabilizing steps that plagued earlier policy gradient methods such as
REINFORCE.

PPO is well-suited to the Central Brain routing task because the state space is
continuous and high-dimensional (25 dimensions) and the training process must
remain stable over a large number of steps. An entropy bonus term with coefficient
$ent\_coef = 0.01$ is added to the objective to encourage exploration throughout
training, preventing the policy from collapsing prematurely to a single strategy.

% ----------------------------------------------------------------
\subsection{Comparison of Reinforcement Learning Techniques}

Table~\ref{tab:rl_comparison} summarizes the properties of the RL techniques
reviewed in this chapter, with reference to their usage in this project.

\begin{table}[h!]
\centering
\caption{Comparison of Reinforcement Learning Techniques}
\label{tab:rl_comparison}
\begin{tabular}{|l|c|c|c|c|}
\hline
\textbf{Property} & \textbf{Tabular Q-Learning} & \textbf{DQN} & \textbf{PPO} & \textbf{Usage in This Project} \\
\hline
Handles Continuous State Space   & No  & Yes & Yes & Required for Central Brain \\
\hline
On-Policy / Off-Policy           & Off & Off & On  & Both modes used \\
\hline
Training Stability               & Low & Medium & High & Critical for 50k-step run \\
\hline
Supports Online Updates          & Yes & No  & No  & Required for Edge Agent \\
\hline
Suitable for Discrete Actions    & Yes & Yes & Yes & All three strategies are discrete \\
\hline
Computational Cost               & Low & High & High & Low for Edge, High for Central \\
\hline
\end{tabular}
\end{table}

\begin{table}[h!]
\centering
\caption{Reinforcement Learning Algorithm Characteristics and Applications}
\label{tab:rl_apps}
\begin{tabular}{|l|l|l|}
\hline
\textbf{Algorithm} & \textbf{Key Characteristics} & \textbf{Application in This Project} \\
\hline
Tabular Q-Learning & Simple, online, low memory & Edge Q-Agent emergency rerouting \\
\hline
Deep Q-Network & Neural network Q-approximation & Not used (Edge is discrete/small) \\
\hline
PPO & Clipped policy gradient, stable & Central Brain strategy selection \\
\hline
Ridge Regression & Regularized linear, online updates & ETA Forecaster travel time estimation \\
\hline
\end{tabular}
\end{table}

% ================================================================
\vspace{1cm}
\noindent \textbf{Chapter Summary:} This chapter reviewed the three categories of
techniques relevant to this project. Traditional routing methods such as A* and
VRP solvers provide deterministic optimal paths but cannot adapt to real-time
events or balance multiple objectives simultaneously. Machine learning methods
improve forecasting accuracy but remain predictive rather than decision-making.
Reinforcement learning, specifically PPO for macro-level strategy selection and
Q-Learning for micro-level emergency rerouting, addresses both the adaptability
and decision-making requirements of the problem. The next chapter describes how
the simulation environment generates data and how that data is transformed into
state representations for the AI models.
